\documentclass[../Main.tex]{subfiles}

\begin{document}

\section{Schlussfolgerungen und weiteres Vorgehen} \label{sec:Schlussfolgerungen}

also damit sind wir zum schluss unsserer arbeit gekommen. in diesem kapitel möchten wir unsere Erkenntnisfrage versuchen zu beantworten und dabei limitation die wir hatten erwähnen und weitere ansätze zur beantwortung unserer Erkenntnisfrage liefern und verbesserungsvorschläge für unseren prototyp machen.

\subsection{Beantwortung der Erkenntnisfrage} \label{sec:Beantwortung}

die erkenntnisfrage war ja: ...

für das gegnerverhalten haben v2 und v3 deutliche hinweise darauf, dass sie so wirken wie sie wirken. und bei v1 und v4 konten wir zwar hinweise finden aber noch nicht so konkrete, wir würden beide fragen weiter verfolgen, da unsere daten mindestens eine gross genug hinweis gibt.

für gegnerplatzierung hat besonders v8 starke hinweise darauf geliefert, dass es was damit zu tun hat. v5-7 aber auch, wenn auch nicht zu stark. wir fürden von denen ich weiss nicht auswendig was 5-7 sind abe einige davon vervolgen was wir dabei genauer anschauen würden ist...
\begin{itemize}[leftmargin=1.5cm]
    \item \textcolor{red}{Erkenntnisfrage wörtlich aufgreifen und explizit beantworten, auf Grundlage von Recherche, Analyse, Prototyp und Playtest (Vertrag C.5).}
    \item \textcolor{red}{Gliederung entlang der zwei Komponenten (Gegnerverhalten, Gegnerplatzierung) und der drei Qualitäten (Spielfluss, Schwierigkeit, Spielgefühl).}
    \item \textcolor{red}{Vorläufig formulieren (\enquote{Für unseren Prototyp zeigen die Playtests, dass \ldots}): Eine abschliessende Antwort ist mit fünf Testpersonen nicht möglich, die Erkenntnisfrage muss laut Vertrag aber beantwortet werden. Die Grenzen folgen direkt in Abschnitt~\ref{sec:Limitationen}.}
    \item \textcolor{red}{Hier auch die 3--5 wichtigsten Erkenntnisse einbauen (früher eigener Abschnitt), ohne Details aus Kapitel~\ref{sec:Evaluation} zu wiederholen.}
    \item \textcolor{red}{Am Schluss den Bogen zum Einstieg in Kapitel~1 schliessen.}
\end{itemize}

\subsection{Limitationen} \label{sec:Limitationen}

Die Ergebnisse dieser Arbeit unterliegen mehreren Einschränkungen.\footnote{Mit Claude formuliert, siehe KI-Deklaration (Abschnitt~\ref{sec:KIDeklaration}).}

Bei den Playtests ist zunächst die kleine Stichprobe zu nennen: Mit fünf Testpersonen lassen sich nur Tendenzen erkennen (Abschnitt~\ref{sec:Auswertung}). Die Urteile in Kapitel~\ref{sec:Evaluation} sind deshalb als Hinweise und nicht als Bestätigungen zu verstehen; wo wir keine Hinweise gefunden haben, bedeutet das nicht, dass das Gegenteil zutrifft. Zudem kannten wir die meisten Testpersonen persönlich oder über Bekannte, was trotz unserer Bitte um ehrliche Antworten zu wohlwollenderen Einschätzungen geführt haben kann (Abschnitt~\ref{sec:AuswahlTP}). Zwei Playtests fanden über Discord statt, wo wir nur Bildschirm und Stimme der Testpersonen wahrnehmen konnten, nicht aber Mimik und Körpersprache (Abschnitt~\ref{sec:Beobachtungsraster}).

Auch das Post-Questionnaire hat Schwächen. Q5 enthält zwei Aussagen auf einmal (überschaubare Anfangsposition und danach unvorhersehbareres Verhalten), sodass unklar bleibt, auf welchen Teil sich eine Bewertung bezieht. Q6 spricht von \emph{dem} engen und \emph{dem} offenen Abschnitt, obwohl der Prototyp mehrere Narrow- und Prospect-Räume enthält (Tabelle~\ref{tab:Raeume}); die Testpersonen könnten dabei unterschiedliche Räume im Kopf gehabt haben.

Bei der Referenzspielanalyse wurde jedes Spiel nur von einer Person analysiert. Die Struktur des MDA-Fragebogens ist zwar einheitlich, Einschätzungen wie \enquote{befriedigend} bleiben aber subjektiv und personenabhängig (Abschnitt~\ref{sec:Analysefragebogen}). Zudem konnten wir die Funktionsweise der Referenzspiele nur aus dem beobachteten Verhalten rekonstruieren, da wir keinen Einblick in deren Quellcode hatten.

\begin{itemize}[leftmargin=1.5cm]
    \item \textcolor{red}{Diesen Abschnitt habe ich aus euren Angaben geschrieben. Falls gewünscht noch ergänzen: mündliche Erklärung der Steuerung vor dem Spielen (statt Selbstentdecken), Gespräche in entspannten Phasen, Leveldesign-Quellen nicht genrespezifisch, Zeitmangel bei der Auswahl der Recherchepunkte.}
\end{itemize}

\subsection{Weiteres Vorgehen und Ausblick} \label{sec:Verbesserungen}

\begin{itemize}[leftmargin=1.5cm]
    \item \textcolor{red}{Optional, als Verbesserungen ausserhalb der Erkenntnisfrage (aus dem gestrichenen Abschnitt \enquote{Weitere Beobachtungen}): Signal am Ende eines Encounters bzw.\ bei offener Tür (TP2, TP4, B10), Combo-Abbruch, Tutorial statt mündlicher Erklärung, Steuerung von Dash und Sprung (TP3, B10; TP2, Q13), Reichweite (TP5, B10), Heilungsdauer (TP4, B10).}
\end{itemize}

\begin{itemize}[leftmargin=1.5cm]
    \item \textcolor{red}{\textbf{Weiteres Vorgehen:} konkrete Verbesserungen am Prototyp, die wir aus den Playtests ableiten (Vertrag C.4: \enquote{Daraus werden konkrete Verbesserungen abgeleitet, die nicht mehr umgesetzt werden müssen}). Das Wort \enquote{Verbesserungen} im ersten Satz verwenden. Jede Verbesserung mit Beleg aus Kapitel~\ref{sec:Evaluation}, z.\,B. akustische Angriffsankündigung, Signal beim Öffnen der Tür, abbrechbare Combos, Tutorial.}
    \item \textcolor{red}{\textbf{Ausblick:} 2--3 Sätze, was man über den Prototyp hinaus untersuchen könnte, z.\,B. mehr Testpersonen, zweiter Playtest nach den Verbesserungen, Spawn-Art als eigene Variable, nicht getestete Prinzipien aus Kapitel~\ref{sec:Leveldesign} (Höhen und Flanken).}
\end{itemize}

\end{document}
